% !TEX root = main.tex 
\section{Introducción} 
El método que se presenta a continuación es una mejora del tradicional que se presenta en el curso de Física Atómica y Molecular, el cual suele estar basado en tablas y contabilidad de posibilidades según estas se le ocurran al estudiante. El método que presentamos reduce el esfuerzo de contabilidad gracias a la combinatoria y simetrías de los términos, sistematizando el proceso y siendo fácilmente extensible a un número arbitrario de electrones.\par
El procedimiento consiste en centrarse en el número cuántico $M_{L}$ (o $M_{J}$ para el método j-j) e ir obteniendo los términos de acuerdo al número de combinaciones posibles de los electrones para dar ese número cuántico. Conviene empezar desde el máximo valor posible de $M_{L}$ (o $M_{J}$) e ir reduciendo su valor hasta llegar a $M_{L}=0$ (o $M_{J}=0$) (hasta ese valor será suficiente gracias a las simetrías). Esta reducción se puede realizar sistemáticamente gracias a la aplicación del operador $L_{-}$ (o $J_{-}$) como veremos en los ejemplos. Según las distintas combinaciones que surjan, y teniendo en cuenta el spin para el caso de L-S, obtendremos los términos correspondientes, que deberemos ir anotando a medida que los vayamos encontrando, puesto que la siguiente reducción los necesitará en la contabilidad. 
